% GENERATED FILE. Edit canonical lessons, then run npm run generate:books.
% canonical-source-hash: fnv1a64:d202c22e63d5e60c
% canonical-lessons: MW-C212-sava, MW-C212-dedh, MW-C212-adhai, MW-C212-duno, MW-C212-chhatak

\chapter{A Quarter and a Half}
\label{ch:mw-a-quarter-and-a-half}
\hlchaptermodality{212}

\begin{chapteropening}
\textbf{By the end of this chapter:} \emph{I can ask for one and a quarter, one and a half or double.}

Complete the last lesson of chapter 212: I can ask for one and a quarter, one and a half or double.
\end{chapteropening}

\section[\texorpdfstring{\mw{सवा} (savā)}{सवा (savā)}]{\mw{सवा} (savā) — plus a quarter}

\label{lesson:MW-C212-sava}



\begin{quote}
Before the new one: say the Marwadi for tight, then the Marwadi for to fit.
\end{quote}

\subsection*{You'll want to know: \mw{सवा}}
\textbf{\mw{सवा}} — \emph{savā} — \textquotedblleft{}plus a quarter\textquotedblright{}.

\begin{quote}
\mw{सवा} \mw{सेर} \mw{आटो} \mw{दो।}
\end{quote}

\begin{quote}
\emph{savā ser āṭo do.}
\end{quote}

\begin{quote}
Give me a seer and a quarter of flour.
\end{quote}

\begin{cousinweb}
\textbf{\mw{सवा}} (\emph{savā}) is inherited from Sanskrit \emph{sapāda}, \textquotedblleft{}increased by a fourth\textquotedblright{}; CDIAL lists Marwari \emph{savā}, \textquotedblleft{}plus a quarter\textquotedblright{} (CDIAL 13134).
\end{cousinweb}

\begin{grammarlens}[title={What you've built}]
One more word for this chapter.
\end{grammarlens}

\subsection*{Your turn}
\begin{itemize}
\raggedright
  \item \emph{Say it:} \emph{savā}
  \item \emph{Say it:} \emph{savā}, once more
  \item \emph{Say it:} \emph{savā ser āṭo do.} — Give me a seer and a quarter of flour.
  \item \emph{Recall:} say the Marwadi for tight, then the Marwadi for to fit, then say \emph{savā} again
\end{itemize}

\begin{tcolorbox}[breakable,colback=teal!4,colframe=teal!35!black,title={Before you move on}]
What does \mw{सवा} mean? (\textquotedblleft{}Plus a quarter\textquotedblright{}.) Say it once more.
\end{tcolorbox}
\section[\texorpdfstring{\mw{डेढ} (ḍeḍh)}{डेढ (ḍeḍh)}]{\mw{डेढ} (ḍeḍh) — one and a half}

\label{lesson:MW-C212-dedh}



\begin{quote}
Before the new one: say the Marwadi for to fit, then the Marwadi for plus a quarter.
\end{quote}

\subsection*{You'll want to know: \mw{डेढ}}
\textbf{\mw{डेढ}} — \emph{ḍeḍh} — \textquotedblleft{}one and a half\textquotedblright{}.

\begin{quote}
\mw{डेढ} \mw{किलो} \mw{चावळ} \mw{दो।}
\end{quote}

\begin{quote}
\emph{ḍeḍh kilo chāvaḷ do.}
\end{quote}

\begin{quote}
Give me one and a half kilos of rice.
\end{quote}

\begin{cousinweb}
\textbf{\mw{डेढ}} (\emph{ḍeḍh}) is inherited from Sanskrit \emph{dvyardha}, \textquotedblleft{}one and a half\textquotedblright{}. CDIAL gives no Marwari form in this entry, but lists Hindi \emph{ḍeṛh} and Sindhi \emph{ḍeḍhu}, \textquotedblleft{}one and a half\textquotedblright{} (CDIAL 6698).
\end{cousinweb}

\begin{grammarlens}[title={What you've built}]
One more word for this chapter.
\end{grammarlens}

\subsection*{Your turn}
\begin{itemize}
\raggedright
  \item \emph{Say it:} \emph{ḍeḍh}
  \item \emph{Say it:} \emph{ḍeḍh}, once more
  \item \emph{Say it:} \emph{ḍeḍh kilo chāvaḷ do.} — Give me one and a half kilos of rice.
  \item \emph{Recall:} say the Marwadi for to fit, then the Marwadi for plus a quarter, then say \emph{ḍeḍh} again
\end{itemize}

\begin{tcolorbox}[breakable,colback=teal!4,colframe=teal!35!black,title={Before you move on}]
What does \mw{डेढ} mean? (\textquotedblleft{}One and a half\textquotedblright{}.) Say it once more.
\end{tcolorbox}
\section[\texorpdfstring{\mw{अढाई} (aḍhāī)}{अढाई (aḍhāī)}]{\mw{अढाई} (aḍhāī) — two and a half}

\label{lesson:MW-C212-adhai}



\begin{quote}
Before the new one: say the Marwadi for plus a quarter, then the Marwadi for one and a half.
\end{quote}

\subsection*{You'll want to know: \mw{अढाई}}
\textbf{\mw{अढाई}} — \emph{aḍhāī} — \textquotedblleft{}two and a half\textquotedblright{}.

\begin{quote}
\mw{ये} \mw{अढाई} \mw{सो} \mw{रुपिया} \mw{है।}
\end{quote}

\begin{quote}
\emph{ye aḍhāī so rupiyā hai.}
\end{quote}

\begin{quote}
This is two hundred and fifty rupees.
\end{quote}

\begin{cousinweb}
\textbf{\mw{अढाई}} (\emph{aḍhāī}) is inherited from Sanskrit \emph{ardhatṛtīya}, \textquotedblleft{}two and a half\textquotedblright{}. CDIAL gives no Marwari form in this entry, but lists Sindhi \emph{aḍhāī} and Hindi \emph{aṛhāī}, \textquotedblleft{}two and a half\textquotedblright{} (CDIAL 651).
\end{cousinweb}

\begin{grammarlens}[title={What you've built}]
One more word for this chapter.
\end{grammarlens}

\subsection*{Your turn}
\begin{itemize}
\raggedright
  \item \emph{Say it:} \emph{aḍhāī}
  \item \emph{Say it:} \emph{aḍhāī}, once more
  \item \emph{Say it:} \emph{ye aḍhāī so rupiyā hai.} — This is two hundred and fifty rupees.
  \item \emph{Recall:} say the Marwadi for plus a quarter, then the Marwadi for one and a half, then say \emph{aḍhāī} again
\end{itemize}

\begin{tcolorbox}[breakable,colback=teal!4,colframe=teal!35!black,title={Before you move on}]
What does \mw{अढाई} mean? (\textquotedblleft{}Two and a half\textquotedblright{}.) Say it once more.
\end{tcolorbox}
\section[\texorpdfstring{\mw{दूणो} (dūṇo)}{दूणो (dūṇo)}]{\mw{दूणो} (dūṇo) — double; twice as much}

\label{lesson:MW-C212-duno}



\begin{quote}
Before the new one: say the Marwadi for one and a half, then the Marwadi for two and a half.
\end{quote}

\subsection*{You'll want to know: \mw{दूणो}}
\textbf{\mw{दूणो}} — \emph{dūṇo} — \textquotedblleft{}double; twice as much\textquotedblright{}.

\begin{quote}
\mw{आज} \mw{भाव} \mw{दूणो} \mw{है।}
\end{quote}

\begin{quote}
\emph{āj bhāv dūṇo hai.}
\end{quote}

\begin{quote}
Today the price is double.
\end{quote}

\begin{cousinweb}
\textbf{\mw{दूणो}} (\emph{dūṇo}) is inherited from a reconstructed \emph{*duguṇa}, \textquotedblleft{}double\textquotedblright{}. CDIAL gives no Marwari form in this entry, but lists Hindi \emph{dūnā}, \textquotedblleft{}double\textquotedblright{} (CDIAL 6390).
\end{cousinweb}

\begin{grammarlens}[title={What you've built}]
One more word for this chapter.
\end{grammarlens}

\subsection*{Your turn}
\begin{itemize}
\raggedright
  \item \emph{Say it:} \emph{dūṇo}
  \item \emph{Say it:} \emph{dūṇo}, once more
  \item \emph{Say it:} \emph{āj bhāv dūṇo hai.} — Today the price is double.
  \item \emph{Recall:} say the Marwadi for one and a half, then the Marwadi for two and a half, then say \emph{dūṇo} again
\end{itemize}

\begin{tcolorbox}[breakable,colback=teal!4,colframe=teal!35!black,title={Before you move on}]
What does \mw{दूणो} mean? (\textquotedblleft{}Double; twice as much\textquotedblright{}.) Say it once more.
\end{tcolorbox}
\section[\texorpdfstring{\mw{छटांक} (chhaṭā̃k)}{छटांक (chhaṭā̃k)}]{\mw{छटांक} (chhaṭā̃k) — a chhatak (a sixteenth of a seer)}

\label{lesson:MW-C212-chhatak}



\begin{quote}
Before the new one: say the Marwadi for two and a half, then the Marwadi for double.
\end{quote}

\subsection*{You'll want to know: \mw{छटांक}}
\textbf{\mw{छटांक}} — \emph{chhaṭā̃k} — \textquotedblleft{}a chhatak (a sixteenth of a seer)\textquotedblright{}.

\begin{quote}
\mw{एक} \mw{छटांक} \mw{जीरो} \mw{दो।}
\end{quote}

\begin{quote}
\emph{ek chhaṭā̃k jīro do.}
\end{quote}

\begin{quote}
Give me a chhatak of cumin.
\end{quote}

A feminine noun.

\begin{cousinweb}
\textbf{\mw{छटांक}} (\emph{chhaṭā̃k}) is inherited from a reconstructed \emph{*ṣaṭṭaṅka}, \textquotedblleft{}weight of six ṭaṅkas\textquotedblright{}. CDIAL gives no Marwari form in this entry, but lists Hindi \emph{chaṭā̃k}, \textquotedblleft{}a sixteenth of a ser\textquotedblright{} (CDIAL 12785).
\end{cousinweb}

\begin{grammarlens}[title={What you've built}]
One more word for this chapter.
\end{grammarlens}

\subsection*{Your turn}
\begin{itemize}
\raggedright
  \item \emph{Say it:} \emph{chhaṭā̃k}
  \item \emph{Say it:} \emph{chhaṭā̃k}, once more
  \item \emph{Say it:} \emph{ek chhaṭā̃k jīro do.} — Give me a chhatak of cumin.
  \item \emph{Recall:} say the Marwadi for two and a half, then the Marwadi for double, then say \emph{chhaṭā̃k} again
\end{itemize}

\begin{tcolorbox}[breakable,colback=teal!4,colframe=teal!35!black,title={Before you move on}]
What does \mw{छटांक} mean? (\textquotedblleft{}A chhatak (a sixteenth of a seer)\textquotedblright{}.) Say it once more.
\end{tcolorbox}
